% Claim proof file for row claim_3_2 -- "AcyclicIffTopologicalOrder".
% Auto-generated stub by scaffold/solve_chapter.py at row start. A manager
% agent dedicated to writing the TeX proof fills in the proof body below.
%
% This file is a sibling of `main.tex` inside `<subsection>/tex/`. The
% `subfiles` package lets it render both standalone and as part of
% `main.tex` via `\subfile{...}`.
%
% The statement is restated above the proof so the file is self-contained
% when rendered on its own. The proof must:
%   - Stay close to the lecture notes' proof if one exists (always search
%     the LN first for a `\begin{proof}` block immediately following the
%     claim; copy / lightly edit when present).
%   - Otherwise use the LN paradigm and cite prior claims by their `ref`.
%   - Be self-contained enough that a mathematician can follow it without
%     re-reading the LN.
%   - Have no `sorry`, `omitted`, or "obvious" shortcuts.

\documentclass[main]{subfiles}

\begin{document}
% ref: claim_3_2 (proof, with statement restated for readability)
% title: AcyclicIffTopologicalOrder
\def\rowref{claim\_3\_2}
\phantomsection\label{claim_3_2}

% --- statement (restated) ---------------------------------------------------
\begin{Lem}[AcyclicIffTopologicalOrder]
        A CDMG  $G=(J,V,E,L)$ is acyclic if and only if it has a topological order.
\end{Lem}

% --- proof ------------------------------------------------------------------
% Proof body lifted verbatim from the LN (lecture_notes/graphs.tex, L228-244),
% with only the outer `\Claude{...}` annotation wrapper stripped.
\begin{proof}
($\Leftarrow$)
Suppose $<$ is a topological order of $G$, i.e.\ $v \in \Pa^G(w) \implies v < w$.
If there were a non-trivial directed walk $v = v_0 \tuh v_1 \tuh \cdots \tuh v_n = v$ ($n \ge 1$),
then $v_0 < v_1 < \cdots < v_n = v_0$, a contradiction.

\medskip\noindent($\Rightarrow$)
Suppose $G$ is acyclic. We construct a topological order by repeatedly selecting a node with no parents in the remaining graph.
Define a sequence $v_1, \dots, v_K$ ($K = |J \cup V|$) inductively: given $v_1,\dots,v_{i-1}$, let $R_i := (J \cup V) \sm \{v_1,\dots,v_{i-1}\}$
and $G_i$ the subgraph of $G$ induced on $R_i$.
Since $G_i$ is acyclic (as an induced subgraph of an acyclic graph) and finite, it has a node $v_i$ with $\Pa^{G_i}(v_i) = \emptyset$
(otherwise, following parent edges would produce a directed cycle).
Choose any such $v_i$.

This gives a total order $v_1 < v_2 < \cdots < v_K$.
If $v = v_j \in \Pa^G(w = v_i)$, then when $w$ was selected, $v$ was still in $R_i$, so $v$ would have been a parent of $w$ in $G_i$, contradicting $\Pa^{G_i}(v_i) = \emptyset$---unless $v$ was already selected, i.e.\ $j < i$. Hence $v < w$.
\end{proof}

\end{document}
